\section{A 227-state port by constrained source reassociation}
\label{sec:source227}

The new candidate has 227 working states, a binary tape alphabet, and one
designated halting state.  This section uses the inherited numerical state
labels, with start state 95, dispatcher 96, and halt label 254.  A separate
renumbering gives the usual consecutive labels.  The construction begins
with the 229-state result of the loop rotation and the state-90/72 fold.
Its additional stages are
\[
 T_{229}\ \longrightarrow\ B_{229}\ \longrightarrow\ B_{228}
 \ \longrightarrow\ T_{227}.
\]
The first arrow changes the source parenthesization and its binary decoder
without changing the number of working states.  The next arrow is a
context-sensitive fold of state 218.  The final arrow is the delayed-erasure
fold of state 169 into state 80.  The intermediate $B_{228}$ is distinct
from the direct 228-state delayed-erasure port obtained from $T_{229}$.

\subsection{Exact scope of the source change}

The source before and after this construction has 5,121 instruction
occurrences.  Numbering these in their common left-to-right order, every
instruction, register operand, and guarded successor is identical.  There
are 6,656 guarded cases, counting the zero and positive cases of each
decrement separately.  Although 4,916 leaf paths change, the ordinal
guarded operational graph does not change.

The source directed acyclic graph has 212 distinct nodes before the
reconstruction and 211 afterwards.  Its maximum leaf depth falls from 40
to 37.  Its longest run of zero bits is six in both the calibrated input
and the final result.  Intermediate source rearrangements can have a
seven-zero run and are not accepted as final binary encodings.

The reconstruction records seven source stages.  The first three transfer
the exact common chunks of stages 1, 5, and 9 of the earlier source103
boundary-fold construction.  The fourth and fifth move the terminal
break through two closed chunks.  The sixth repairs an entry context
specific to K230.  The final stage reparses an ordered sequence of 49
closed atoms.

\begin{center}
\begin{tabular}{rlll}
Stage & Operation & Chunk sizes & Occurrences\\\hline
1 & left reassociation & $19,7,26$ & 10\\
2 & left reassociation & $189,212,651$ & 1\\
3 & left reassociation & $2,82,80$ & 1\\
4 & right reassociation ending in break 1 & $81,4853,1$ & 1\\
5 & right reassociation ending in break 1 & $527,4326,1$ & 1\\
6 & right reassociation at states 502/501 & $52,54,189$ & 1\\
7 & ordered 49-atom reassociation & 4,326 leaves in total & 1
\end{tabular}
\end{center}

Here a left reassociation replaces $(A,(B,C))$ by $((A,B),C)$, and a
right reassociation makes the inverse change.  These names describe
parenthesization, not movement of the Turing-machine head.

\subsection{An elementary source-equivalence proof}

For an instruction tree $U$, let
$\partial U=(b_U,c_U,h_U)$ bound its escaping break levels, escaping
continue levels, and the presence of a halt instruction.  On leaves,
\[
\begin{aligned}
 \partial\mathrm{CI}(r)&=(0,0,\mathrm{false}),&
 \partial\mathrm{CD}(r)&=(1,0,\mathrm{false}),\\
 \partial\mathrm{CB}(k)&=(k,0,\mathrm{false}),&
 \partial\mathrm{CC}(k)&=(0,k,\mathrm{false}),\\
 \partial\mathrm{CH}&=(0,0,\mathrm{true}).
\end{aligned}
\]
Writing $x^+=\max(x,0)$, the recursive clause is
\[
 \partial(A,B)=
 \bigl(\max((b_A-1)^+,b_B),
       \max(c_A,(c_B-1)^+),h_A\lor h_B\bigr).
\]
An atom with signature $(0,0,\mathrm{false})$ is closed: its nonlocal
control instructions resolve inside the atom.  Any parenthesization of
an ordered list of such atoms preserves their register actions and
control destinations.  A fall-through still reaches the next ordinal
instruction.  A divergent execution inside an atom is preserved as
well, so this argument does not assume that any atom terminates.

Stages 1--3 and 6 have three closed chunks.  Every one of the 49 atoms in
stage 7 is also closed.  In stages 4 and 5 the first two chunks are
closed and the last chunk is the single instruction $\mathrm{CB}(1)$.
The latter instruction changes its suffix from $1$ to $11$ under the
right reassociation.  No zero is inserted on its ancestry, so the
nearest enclosing zero edge and hence its break destination are
unchanged.  These observations establish each source stage directly.

The executable reconstruction additionally computes the complete guarded
source graph twice at every stage.  One resolver uses binary successor
and path cutting.  It checks that every unused suffix after resolution
contains only zeros, the required no-stray-bit condition.  The other
resolver independently assigns ordinal intervals to trie nodes and
walks outward through ancestors to resolve breaks and continues.  Both
resolvers agree with each other and with the input graph, and the
initial ordinal leaf is zero throughout.

Consequently, corresponding source configurations with any nine
natural-valued registers make identical register updates and select
identical ordinal successors.  Induction proves equality of every finite
source run.  The same step correspondence preserves infinite runs and
halting.  This is a proof for all register values obtained from a finite
description of the operational graph.

\subsection{The new layout search}

The successful search treats source parenthesization as an ordered
binary-tree optimization subject to decoder constraints.  The affected
component is rooted at inherited state 528, at source prefix
\texttt{11101001001110}.  Its 49 atoms contain 4,326 instruction
occurrences.  Its 48 internal decoder labels are reused.  The atoms,
their order, and their internal source trees are fixed.

For an entry context, retain the last two path bits and the current
length of the trailing zero run.  Appending a left or right edge updates
this finite context.  The discovery filter rejects a zero run exceeding
six, rejects changes to the inherited two-bit entry tags of the selected
nontrivial decoder nodes, and requires every source entry to state 179
to have a preceding zero.  States serving the startup and backward-read
helper paths are placed inside fixed atoms rather than freely reparsed.
These filters guide discovery; the final symbolic runtime proof remains
the correctness check for arbitrary counters and spares.

Let $A_i,\ldots,A_{j-1}$ be an interval of atoms and let $C$ be its set of
possible incoming contexts.  If $L(C)$ and $R(C)$ append a zero and a one,
respectively, and $W_{ij}$ is the number of instruction occurrences in
that interval, the dynamic program uses
\[
 D(i,j,C)=\min_{i<k<j}
 \bigl(D(i,k,L(C))+D(k,j,R(C))+W_{ij}\bigr).
\]
For a single atom the cost is zero when all its fixed paths satisfy the
constraints, and is infinite otherwise.  The objective is the sum of
atom entry depths weighted by atom sizes.  For this component it falls
from 27,716 to 26,568.  The result is optimal for this fixed atom sequence,
objective, and context filter; no global minimality claim is made for
the source or the Turing machine.

The deliverable does not need to rerun this search.  The construction
record contains the literal old and new binary patterns, numbers every
atom from 0 through 48 in ordinal order, pins each atom by its source
hash, and records all 48 reused state labels.  The replay verifies that
neither pattern drops, duplicates, or permutes an atom.  Thus optimizer
implementation or tie-breaking choices are outside the replay's trust
requirements.

\subsection{Decoder and helper repairs}

Recompiling the new source into $B_{229}$ changes 62 literal transition
rows.  The complete edit set and every old-row precondition are in the
construction record.  Three helper adjustments deserve explicit mention.
First, state 61 acquires a changed forward decoding role, so the backward
use at $120/1$ is redirected to the already existing state 202, whose
required branch has the same behavior.  Second, the terminal break's
longer one suffix is handled by
\[
 138/1=0\ L\ 120.
\]
Third, state 529 is reused as the internal helper of $218/0$:
\[
 218/0=0\ L\ 529,\qquad
 529/0=0\ R\ 16,\qquad 529/1=1\ R\ 99.
\]
This restores a backward branch whose former decoder continuation was
changed by source reassociation.  The basis therefore still has 229
states.  The purpose of these edits is to obtain the useful entry
contexts without introducing an additional helper state.

\subsection{The boundary fold from 229 to 228 states}

In $B_{229}$ the only incoming transitions to state 179 are $236/0$ and
$501/0$, both writing zero and moving right.  State 218 has three
incoming ports: $187/1$, $220/1$, and $237/0$, all moving right.  The
sole incoming transition to helper 529 is $218/0$, a left move.

Delete state 218.  Redirect each old right-moving call of 218 to state
529 if that call writes zero, or to state 179 if it writes one.  Make
the following three row replacements:
\[
 179/1=1\ L\ 215,\qquad
 529/0=0\ R\ 23,\qquad 529/1=1\ R\ 70.
\]
All other ordinary transitions are unchanged.  The finite certificate
checks these six block identities, where $t$ is the preceding tape bit,
$b$ is the scanned bit, and every endpoint has head displacement $+1$:
\begin{center}
\begin{tabular}{rrrrrrr}
Old state & $t$ & $b$ & New state & Old steps & New steps & Endpoint\\\hline
179&0&0&179&1&1&105\\
179&0&1&179&1&5&105\\
218&0&0&529&3&1&23\\
218&0&1&529&1&1&70\\
218&1&0&179&3&1&105\\
218&1&1&179&1&3&70
\end{tabular}
\end{center}
Each block reads only the preceding and scanned cells and restores both
exactly.  All other tape cells are arbitrary and untouched.  Old helper
529 is internal to the three-step $218/0$ block and is not a separate
checkpoint.  The incoming-port check establishes the preceding-bit
premises whenever an exceptional checkpoint is entered.  The start state
is ordinary, and each paired block has positive length in both machines.
Composing the blocks proves a whole-tape simulation from start state 95
on every binary tape, preserving both finite halting and infinite runs.

\subsection{The final fold and the blank start}

The delayed-erasure fold then deletes state 169 and makes
\[
 147/0=0\ R\ 80,\qquad
 80/0=0\ L\ 106,\qquad
 91/1=1\ L\ 167.
\]
The resulting table has 227 working states and 454 transition rows.
The exact table, source, and stage hashes are checked by the replay.
The independent unbounded binary-runtime proof is applied directly to
this final table and the reconstructed source; it is not inferred from
a bounded instruction test.

The replay also executes the complete blank startup of every displayed
stage.  The input $T_{229}$, source-recompiled $B_{229}$, and split
$B_{228}$ each reach dispatcher 96 after 409,691 transitions.  The final
$T_{227}$ reaches the same dispatcher after 409,835 transitions.  Every
endpoint has head position $-5$ and exactly the same 351 marked cells;
all other cells are zero.  Comparing the full support of the tapes
establishes equality of the complete initialization frames.

The source equivalence, startup equality, and the final unbounded runtime
certificate are separate obligations.  Together with the documented
theorem for the pinned K230 ancestor, they support transfer of its
blank-start halting characterization.  The new finite certificates and
their mathematical composition argument are not presented as a newly
checked Rocq theorem for the 227-state table.
